\documentclass[10pt,a4paper]{amsart}
\usepackage[margin=19mm]{geometry}
\usepackage{amsmath,amssymb,mathtools}
\usepackage[scr=rsfs,cal=euler]{mathalfa}
\usepackage{xfrac}
\usepackage[unicode,hidelinks,bookmarks=false]{hyperref}
\newcommand{\ie}{i.e.\ }
\newcommand{\cf}{cf.\ }
\newcommand{\resp}{respectively\ }
\newcommand{\Subsection}{\subsection}
\providecommand{\nobreakdash}{\nobreak\hbox{-}}
\setlength{\emergencystretch}{2em}
\multlinegap0pt
\usepackage{bm}
\usepackage{bbm}
\usepackage{dsfont}
\usepackage{xspace}
\usepackage{cancel}
\usepackage{mathtools}
\usepackage{amsthm}
\makeatletter
\@ifundefined{definition}{\newtheorem{definition}{Definition}}{}
\@ifundefined{proposition}{\newtheorem{proposition}[definition]{Proposition}}{}
\@ifundefined{theorem}{\newtheorem{theorem}[definition]{Theorem}}{}
\makeatother
\newcommand{\M}{\mathcal{M}}
\newcommand{\N}{\mathbb{N}}
\title[Law of large numbers for non-linear traces of the Choquet ty]{Law of large numbers for non-linear traces of the Choquet type on finite factors}
\author{Masaru Nagisa, Yasuo Watatani}
\date{Source: arXiv:2510.22280v1 (CC BY 4.0)}
\begin{document}
\maketitle

\noindent\emph{Source and licence.} Law of large numbers for non-linear traces of the Choquet type on finite factors. Version arXiv:2510.22280v1 (CC BY 4.0), distributed under the Creative Commons Attribution 4.0 International licence, as recorded in the arXiv licence metadata for this exact version and shown on the abstract page https://arxiv.org/abs/2510.22280v1. Full rights notice: licenses/arxiv-2510.22280-CC-BY-4.0.txt.\par\smallskip

\noindent\textit{Editorial scope.} Counted unit: the Theorem whose source label is \texttt{psudo law of large numbers} in the article (source0.tex lines 1130--1155, source label \texttt{psudo law of large numbers}), delivered together with the article's own complete original proof of it (source0.tex lines 1156--1220, one \texttt{proof} environment of 65 lines). No mathematics is written, repaired or re-derived: statement and proof are the article's own text line for line. Required context retained uncounted: prop:tranlatable (lines 717--742); superadditivity (lines 919--942). Every definition, display and label of those lines is retained unchanged. Omitted from this package and not redistributed: 1--716, 743--918, 943--1129, 1221--2103. Nothing inside a retained excerpt is omitted or altered: every retained body is delivered exactly as the article's lines are written, so no omission receipt of any kind accompanies this package. Citations: the retained excerpts carry no citation command at all, so no bibliography entry of the article is transcribed and no reference is redistributed. Reference adaptation: none was needed and none was made; every reference inside the delivered unit resolves inside it, so no unresolved reference or citation is produced. Printed slips kept literal: no printed word, symbol, hypothesis or number of the counted statement or of its proof was altered. Layout and class: the article's own class is \texttt{amsart} with packages including graphicx, amsmath, amssymb, amsfonts, amsthm, color; the wrapper is the lane's pinned \texttt{amsart} editorial wrapper at 10pt on A4, which restates the theorem-like environments the retained text needs and adds the article's own macro definitions that the retained text needs (\texttt{\textbackslash M}, \texttt{\textbackslash N}), with extra wrapper packages (bm, bbm, dsfont, xspace, cancel, mathtools). The delivered numbering is the unit's own. Assets: no retained span contains a figure environment, a graphics-inclusion command, a PDF-page inclusion, a table or a TikZ picture; no image file is copied into this package, the package declares no asset dependency and no image file is redistributed with it. Licence: version arXiv:2510.22280v1 of this article is distributed under the Creative Commons Attribution 4.0 International licence, as recorded in the arXiv licence metadata for this exact version and shown on the abstract page https://arxiv.org/abs/2510.22280v1. A full rights notice is delivered at licenses/arxiv-2510.22280-CC-BY-4.0.txt. Characters: every retained excerpt is pure ASCII, so the delivered engine, XeLaTeX with its default Latin Modern text font, reports no missing character.
\begin{proposition} 
\label{prop:tranlatable}
Let $\M$ be a factor of type ${\rm II}_1$ with a normalized trace $\tau$.  
Let $\alpha: [0,1]  \rightarrow [0, 1]$ be a 
monotone increasing continuous function with $\alpha(0) = 0$ and $\alpha(1) = 1$ 
Then the translatable non-linear trace $\varphi_{\alpha}^{t}$ of the Choquet type on $\M_{s.a.}$ satisfies the following:
\begin{enumerate}
\item[(1)] {\rm (positively homogeneity)} $\varphi_{\alpha}^{t}(ka) = k \varphi_{\alpha}^{t}(a)$ for any 
$a \in  {\M}_{s.a.}$ and any scalar $k \geq 0$. \\
\item[(2)] {\rm (translatability)}  For any $a \in  {\M}_{s.a.}$ and any real number $c$ 
$$
\varphi_{\alpha}^{t}(a + cI) = \varphi_{\alpha}^{t}(a) + c.
$$
\item[(3)] {\rm (monotonicity)} For $a, b \in {\M}_{s.a.}$, if $a \leq b$, then $\varphi_{\alpha}^{t}(a) \leq 
\varphi_{\alpha}^{t}(b). $\\
\item[(4)] {\rm (unitary invariance)} For any $a \in  \M_{s.a.}$ and any unitary $u \in \M$, we have that 
$\varphi_{\alpha}^{t}(uau^*) =  \varphi_{\alpha}^{t}(a)$. \\
\item[(5)] {\rm (triangle inequality)} Suppose that $\alpha$ is concave on $[0,1]$. Then 
for any $a, b \in {\M}_{s.a.}$, 
$$
 \varphi_{\alpha}^{t}(a + b) \leq  \varphi_{\alpha}^{t}(a) + \varphi_{\alpha}^{t}(b).
$$

\item[(6)] {\rm (operator norm continuity)} $\varphi_{\alpha}^{t}$ is operator norm continuous. 
\end{enumerate}
\end{proposition}

\begin{proposition} \label{superadditivity}
Let $\M$ be a factor of type ${\rm II}_1$ with a normalized trace $\tau$.  
Let $\alpha: [0,1]  \rightarrow [0, 1]$ be a 
monotone increasing continuous function with $\alpha(0) = 0$ and $\alpha(1) = 1$ 
\begin{enumerate}
\item[$(1)$] 
If $\alpha$ is concave on $[0,1]$, then $\overline{\alpha}$ is convex on $[0,1]$. 
If $\alpha$ is convex on $[0,1]$, then $\overline{\alpha}$ is concave on $[0,1]$. 
\item[$(2)$] 
If $\alpha$ is convex on $[0,1]$, then for any $a, b \in {\M}_{s.a.}$, 
\[  
\varphi_{\alpha}^{t}(a + b) \geq \varphi_{\alpha}^{t}(a) + \varphi_{\alpha}^{t}(b).
 \]
\item[$(3)$] 
If $\alpha$ is concave on $[0,1]$, then  for any $a \in {\M}_{s.a.}$,
\[ 
\varphi_{\overline{\alpha}}^{t}(a)  \leq \varphi_{\alpha}^{t}(a).
\]
If $\alpha$ is convex on $[0,1]$, then  for any $a \in {\M}_{s.a.}$,
\[
\varphi_{\alpha}^{t}(a) \leq \varphi_{\overline{\alpha}}^{t}(a). 
\]
\end{enumerate}
\end{proposition}

\begin{theorem} \label{psudo law of large numbers}
Let $\M$ be a factor of type ${\rm II}_1$ with a normalized trace $\tau$.  
Let $\alpha: [0,1]  \rightarrow [0, 1]$ be a monotone increasing continuous function with $\alpha(0) = 0$ 
and $\alpha(1) = 1$, and
$\varphi_{\alpha}$ a non-linear trace of the Choquet type associated with $\alpha$.  
Let $(a_n)_n$ be an operator norm bounded sequence $(a_n)_n$ in ${\M}_{s.a.}$. 
Put $s_n = a_1 + a_2 + \dots +a_n$.  
\begin{enumerate}
\item[(1)] Assume that  $\alpha$ is concave. Fix $k \in  \N$. Suppose that there is a constant 
$C(k) \in {\mathbb R}$ such that  
for any mutually different $i_1, i_2, \dots, i_k \in \N$, 
$$
\varphi_{\alpha}({\rm Re} (a_{i_1}a_{i_2} \dots a_{i_k}) ) \leq C(k).  
$$
Then we have that 
\[  \limsup _{n \to \infty}  \varphi_{\alpha}((\frac{s_n}{n})^k) \leq C(k) .\]
\item[(2)] Assume that  $\alpha$ is convex. Fix $k \in  \N$. Suppose that there is a constant 
$\hat{C}(k) \in {\mathbb R}$ such that  
for any mutually different $i_1, i_2, \dots, i_k \in \N$,   
$$
\varphi_{\alpha}({\rm Re} (a_{i_1}a_{i_2} \dots a_{i_k}) ) \geq \hat{C}(k). 
$$
Then we have that 
\[  \liminf _{n \to \infty}  \varphi_{\alpha}((\frac{s_n}{n})^k) \geq \hat{C}(k) .\]
\end{enumerate}
\end{theorem}

\begin{proof}
Since  $(a_n)_n$ be an operator norm bounded sequence $(a_n)_n$, there exists a constant $K > 0$ such that 
$||a_n||\leq K$ for all $n \in {\Bbb N}.$ 

(1)  Fix $k \in {\Bbb N}$. When $k = 1$, the assertion is clear by a triangle inequality in Proposition \ref{prop:tranlatable}. Hence we may assume that $k \geq 2.$ 
Since $n$ goes to $\infty$, we may assume that $n \geq k$.
Put 
$$
D_k(n) := \{i = (i_1,i_2,\dots,i_k) \in \{1,2,\dots,n\}^k \ | \  i_1,i_2,\dots,i_k \text{  are all distinct } \} .
$$
The cardinality $|D_k(n)|$ of $D_k(n)$ is $\dfrac{n!}{(n-k)!}$. 
Since the cardinality of $D_k(n)^c$ is $|D_k(n)^c|=n^k-|D_k(n)|$, we have
\[  \lim_{n\to\infty} \frac{|D_k(n)|}{n^k} = \lim_{n\to\infty}\frac{n}{n}\frac{n-1}{n}\cdots\frac{n-k+1}{n} = 1  
\text{ and } \lim_{n\to\infty} \frac{|D_k(n)^c|}{n^k}= 0.  \]
% Put $g_k(n) := d_k(n) - n^k $.  Then 
% as a polynomial function of $n$, $\deg  g_k(n) = k-1$.  The cardinality of $(D_k(n))^c$ is 
% $\sharp((D_k(n))^c) = n^k -d_k(n) = -g_k(n)$.  

We also have two inequalities as follows:
$$
|| {\rm Re}(a_{i_1}a_{i_2} \dots a_{i_k}) || \leq ||  a_{i_1}a_{i_2} \dots a_{i_k}    || \leq K^k, 
$$
for any $i_1, i_2, \dots, i_k$.  And by assumption, 
$$
\varphi_{\alpha}({\rm Re} (a_{i_1}a_{i_2} \dots a_{i_k})) \leq C(k), 
$$
for any mutually different $i_1, i_2, \dots, i_k$. 

Since $\sum_{i \in D_k(n)} a_{i_1}a_{i_2}\cdots a_{i_k}$ is self-adjoint, so is
$\sum_{i \in D_k(n)^c} a_{i_1}a_{i_2}\cdots a_{i_k}$.
Using a triangle inequality in Proposition \ref{prop:tranlatable} (5), we have that
\begin{align*}
     & \varphi_{\alpha}((\frac{s_n}{n})^k) = \frac{1}{n^k}  \varphi_{\alpha}((a_1 + a_2 + \cdots + a_n)^k)  \\
  = & \frac{1}{n^k}  \varphi_{\alpha}( \sum_{i \in D_k(n)} a_{i_1}a_{i_2} \cdots a_{i_k} 
         +  \sum_{i \in (D_k(n))^c} a_{i_1}a_{i_2} \cdots a_{i_k} )  \\
  \leq &  \frac{1}{n^k}  \varphi_{\alpha}( \sum_{i \in D_k(n)} a_{i_1}a_{i_2} \cdots a_{i_k} )
         +  \frac{1}{n^k} \varphi_{\alpha} (\sum_{i \in (D_k(n))^c} a_{i_1}a_{i_2} \cdots a_{i_k} ).
\end{align*}
Examine each part.  The first part is 
\begin{align*}
& \frac{1}{n^k}  \varphi_{\alpha}( \sum_{i \in D_k(n)} a_{i_1}a_{i_2} \cdots a_{i_k} ) 
= \frac{1}{n^k}  \varphi_{\alpha}( \sum_{i \in D_k(n)} {\rm Re} (a_{i_1}a_{i_2} \cdots a_{i_k}) ) \\
\leq &  \frac{1}{n^k}  ( \sum_{i \in D_k(n)} \varphi_{\alpha}({\rm Re}(a_{i_1}a_{i_2} \cdots a_{i_k}) ))
  \leq  \frac{1}{n^k}  ( \sum_{i \in D_k(n)} C(k) ) \\
= & \frac{|D_k(n)|}{n^k} C(k) \to C(k) \quad (n\to \infty).
%& = \frac{1}{n^k}(n^k + g_k(n))C(k)= (1 + \frac{g_k(n)}{n^k})C(k)
%   \to C(k)  \quad (n \to \infty),
\end{align*}
% since the degree of $g_k(n)$ as a polynomial of $n$ is $k-1$.  
The second part is 
\begin{align*} 
  &  \frac{1}{n^k} \varphi_{\alpha} (\sum_{i \in (D_k(n))^c} a_{i_1}a_{i_2} \cdots a_{i_k} ) 
      \leq \frac{1}{n^k}  ||\sum_{i \in (D_k(n))^c} a_{i_1}a_{i_2} \cdots a_{i_k} || \\
 \leq & \frac{1}{n^k}  \sum_{i \in (D_k(n))^c} ||a_{i_1}a_{i_2} \dots a_{i_k} || 
      \leq  \frac{1}{n^k}  \sum_{i \in (D_k(n))^c} K^k \\
   = &  \frac{|D_k(n)^c|}{n^k} K^k \to 0 \quad (n \to \infty).
\end{align*} 
Therefore 
\[   \limsup _{n \to \infty}  \varphi_{\alpha}((\frac{s_n}{n})^k) 
     \leq \limsup _{n \to \infty} \frac{|D_k(n)|}{n^k}C(k)    
           + \limsup _{n \to \infty} \frac{|D_k(n)^c|}{n^k}K^k = C(k).  
\]

$(2)$  Use the inequality in Proposition \ref{superadditivity} instead of a triangle inequality in Proposition \ref{prop:tranlatable} (5). Then we get the desired assertion in a similar way.  \\
\end{proof}

\end{document}
